% Generated draft: requires mathematical and source review.
\documentclass[11pt]{article}
\usepackage[margin=25mm]{geometry}
\usepackage{fontspec}
\setmainfont{Latin Modern Roman}
\usepackage{amsmath,amssymb,mathtools,unicode-math}
\setmathfont{Latin Modern Math}
\usepackage{enumitem,needspace,xurl,hyperref}
\setlength{\parindent}{0pt}
\setlength{\parskip}{5pt}
\setlength{\emergencystretch}{3em}
\newcommand{\N}{\mathbb{N}}
\newcommand{\Z}{\mathbb{Z}}
\newcommand{\Q}{\mathbb{Q}}
\newcommand{\R}{\mathbb{R}}
\newcommand{\C}{\mathbb{C}}
\newcommand{\F}{\mathbb{F}}
\newcommand{\A}{\mathbb{A}}
\newcommand{\PP}{\mathbb{P}}
\newcommand{\E}{\mathbb{E}}
\newcommand{\eps}{\varepsilon}
\newcommand{\dd}{\,\mathrm{d}}
\begin{document}
\section*{UnsolvedMath ID: 4000019}
\textbf{Source problem code:} catalogue entry
\subsection{Definitions and mathematical statement}
The source asks about an Alexandrov space \((X,d)\) with a positive lower sectional-curvature bound. In the usual lower-bound convention this means a geodesic metric space satisfying the appropriate local triangle comparison with the two-dimensional model of curvature \(K>0\). Dimension, completeness, boundary, and measure assumptions must be specified for the intended theorem; the excerpt does not provide them.

For a chosen Markov kernel \(x\mapsto m_x\) with finite first moments, its Ollivier coarse Ricci curvature at distinct points is
\[
 \kappa(x,y)=1-\frac{W_1(m_x,m_y)}{d(x,y)},\qquad
 W_1(\mu,\nu)=\inf_{\gamma\in\Pi(\mu,\nu)}\int d(u,v)\,\gamma(\dd u,\dd v).
\]
Here \(\Pi(\mu,\nu)\) denotes couplings of the two probability measures. Is there a natural, geometrically specified choice of kernels for which positive Alexandrov sectional curvature forces positive coarse Ricci curvature? Can the desired implication be proved by approximation by manifolds or by a suitable parallel-transport construction within Alexandrov spaces?

The kernel choice is an essential part of this question. Merely asking for the existence of any kernel would not capture a geometric implication: an \(x\)-independent probability kernel has \(W_1(m_x,m_y)=0\) for every pair. A meaningful target must therefore state its locality or other naturalness requirements. If a scale-dependent family \(m_x^{\varepsilon}\) is intended, the question must also specify the allowed scales, whether positivity holds for all or only nearby pairs, whether a uniform positive lower bound is sought, and whether curvature is rescaled as \(\varepsilon\to0\). None is fixed in the supplied statement.

A manifold-approximation proof would require an available approximating class and convergence of the selected kernels and transport estimates; the excerpt asserts no such approximation theorem. Likewise it gives no definition of parallel transport through singularities. The research summary attributes related transport/curvature results to Ohta and others, but says these do not fully prove the direct Ollivier conclusion for natural kernels. The 2008 open label has no recorded later authoritative resolution check; the registry's partial-progress label is a supplied assessment.
\subsection{Short English statement}
Does positive Alexandrov sectional curvature imply positive Ollivier coarse Ricci curvature for a suitably specified natural Markov kernel? The source also asks about manifold approximation and parallel transport as proof methods, while leaving the kernels, scales, and space hypotheses unspecified.
\subsection{Sources}
\begin{itemize}
\item Local registry: \url{data/unsolvedmath/problems.json}, permanent UnsolvedMath ID \(4000019\), source code catalogue entry; supplied statement, background, and research summary, updated 2026-08-06.
\item Yann Ollivier, \emph{Discrete Ricci curvature: Open problems} (2008), Problem S: \url{https://www.yann-ollivier.org/rech/publs/problems_curvmarkov.pdf}. Supplied PDF-text extraction accessed 2026-07-29; no later authoritative resolution check is recorded.
\item The research summary mentions related results of Ohta and others but provides no publication titles, precise curvature definitions, or theorem statements.
\end{itemize}
\end{document}
